\thispagestyle{plain}
\chapter*{Resumen}
% Extensión: unas 340 palabras. Cifras tomadas de los Caps. 4 y 5 (tab:res_supervisados, tab:res_pruebas_vlm,
% sec:res_melanoma, sec:res_fuga, sec:res_hiba, tab:res_gen_cascada, tab:res_gen_pipeline, sec:res_gen_concordancia,
% tab:res_costo_generadores). El abstract repite la misma estructura y las mismas cifras.

El diagnóstico de una lesión de piel a partir de una imagen dermatoscópica exige, además de clasificarla, redactar un
reporte que registre lo observado y la conclusión. Los modelos visión--lenguaje (VLM) prometen resolver las dos tareas
con un mismo modelo, pero no está establecido si clasifican las lesiones tan bien como los clasificadores entrenados
solo con imágenes, ni si los reportes que generan afirman el diagnóstico correcto y describen la lesión de la imagen.
Este trabajo estudia ambas preguntas sobre \emph{HAM10000}, con una partición disjunta por lesión, y bajo la
restricción de usar una sola GPU de consumo.

En la clasificación se comparan cuatro clasificadores supervisados, dos VLM generativos de propósito general y un
modelo visión--lenguaje dermatológico, DermLIP. El mejor supervisado, Swin-Tiny, obtiene un F1 macro de
$0.733 \pm 0.026$ en la prueba. Los VLM generativos quedan cerca del azar, y DermLIP sin ajuste queda entre 0.15 y
0.18 de sensibilidad macro por debajo de los mejores supervisados, aunque con una cabeza entrenada sobre su
representación se acerca a ellos. En melanoma, un umbral fijado en validación lleva la detección de Swin-Tiny de tres a
cuatro de cada cinco casos. La evaluación muestra, además, que una partición por imagen sobreestima la sensibilidad
macro en unas 0.07, y que en otra población, \emph{HIBA}, el clasificador pierde cerca del 40\,\% de ella.

Para la generación se propone un \emph{pipeline} que entrega al generador el diagnóstico de Swin-Tiny y decide, según
la concordancia con un \emph{k}-vecinos sobre la representación de DermLIP, si afirmarlo, presentar dos candidatos o
derivar el caso a un especialista. El reporte transcribe el diagnóstico recibido, de modo que su consistencia
diagnóstica, 0.849, coincide con la exactitud del clasificador, y el ruteo la eleva a 0.870. Entre los casos que el
\emph{pipeline} afirma, el clasificador acierta el 91.5\,\%, aunque los errores que comparten los dos predictores se
siguen afirmando. El generador, Mini-VLM, ocupa como máximo 8.6~GB y escribe un reporte en unos 16 segundos.

El trabajo establece la fidelidad diagnóstica de los reportes con métricas automáticas, pero no su validez clínica,
que requiere la evaluación de especialistas, y sus resultados se limitan a la dermatoscopía.

\keywordses{lesiones de piel, dermatoscopía, modelos visión--lenguaje, clasificación diagnóstica, detección
de melanoma, generación de reportes clínicos, incertidumbre}

\pdfbookmark[1]{Abstract}{abstract}
\chapter*{Abstract}

Diagnosing a skin lesion from a dermoscopic image requires not only classifying it but also writing a report that
records the findings and the conclusion. Vision--language models (VLMs) promise to solve both tasks with a single
model, but it is not established whether they classify lesions as well as classifiers trained on images alone, nor
whether the reports they generate state the correct diagnosis and describe the lesion in the image. This work studies
both questions on \emph{HAM10000}, with a lesion-disjoint split, under the constraint of a single consumer GPU.

For classification, four supervised classifiers, two general-purpose generative VLMs and a dermatology
vision--language model, DermLIP, are compared. The best supervised model, Swin-Tiny, reaches a macro F1 of
$0.733 \pm 0.026$ on the test set. The generative VLMs stay close to chance, and zero-shot DermLIP falls 0.15 to 0.18
of macro sensitivity below the best supervised models, although with a head trained on its representation it comes
close to them. For melanoma, a threshold set on validation raises Swin-Tiny's detection from three to four out of five
cases. The evaluation also shows that an image-level split overestimates macro sensitivity by about 0.07, and that on
another population, \emph{HIBA}, the classifier loses about 40\,\% of it.

For report generation, a pipeline is proposed that passes the Swin-Tiny diagnosis to the generator and decides,
according to its agreement with a \emph{k}-nearest-neighbours vote on the DermLIP representation, whether to assert
it, present two candidates or refer the case to a specialist. The report transcribes the diagnosis it receives, so
its diagnostic consistency, 0.849, matches the classifier's accuracy, and routing raises it to 0.870. Among the cases
the pipeline asserts, the classifier is correct in 91.5\,\%, although the errors shared by both predictors are still
asserted. The generator, Mini-VLM, uses at most 8.6~GB and writes a report in about 16 seconds.

This work establishes the diagnostic fidelity of the reports with automatic metrics, but not their clinical validity,
which requires evaluation by specialists, and its results are limited to dermoscopy.

\keywordsen{skin lesions, dermoscopy, vision--language models, diagnostic classification, melanoma
detection, clinical report generation, uncertainty}
